\section{Rank-controlled polylogarithmic generators}
\label{sec:polylogs}

\subsection{What the harmonic rank does and does not imply}
For the rational weight of Section~\ref{sec:classification}, define
\[
 S_{r,N}(R)=\sum_{0\le n_1<\cdots<n_r\le N}
             \Delta(\bm n)^2\prod_iR(n_i)=\det M_N.
\]
No infinite-sum convergence is needed for this finite identity.
By Proposition~\ref{prop:pencil}, $S_{r,N}$ is a polynomial in
shifted harmonic symbols with rational functions of $N$ as coefficients.

\begin{proposition}[Rank bounds harmonic degree]\label{prop:degree}
If the pencil $\mathscr C(\bm X)$ has generic rank $\rho$, then
\begin{equation}\label{eq:degree-bound}
 \deg_{\bm X}\det(D_N+\mathscr C(\bm X))\le\rho.
\end{equation}
In particular, a rank-one pencil makes $S_{r,N}$ affine in all the
harmonic symbols, even with repeated poles and arbitrarily many particles.
A zero pencil makes $S_{r,N}$ rational in $N$.
\end{proposition}
\begin{proof}
Over $K(N,\bm X)$ a rank-$\rho$ matrix is a product of an $r\times\rho$
and a $\rho\times r$ matrix. The determinant
$\det(D_N+t\mathscr C(\bm X))$ therefore has degree at most $\rho$
in $t$, by multilinearity or by the vanishing of minors of order
larger than $\rho$. The coefficient of $t^k$ is exactly the
homogeneous degree-$k$ part in $\bm X$. This proves the assertion
without any assumption that $D_N$ is invertible.
\end{proof}

The degree can be strictly smaller than $\rho$ because of additional
cancellations involving $D_N$. Neither $\rho$ nor the polynomial degree
proves minimal depth of evaluated periods. These are finite structural
upper bounds.

\subsection{A constructive depth bound on rational grids}
For positive integers $p_1,\ldots,p_\ell$, our multiple-polylogarithm
convention is
\begin{equation}\label{eq:MPL}
 \Li_{p_1,\ldots,p_\ell}(z_1,\ldots,z_\ell)
 =\sum_{n_1>\cdots>n_\ell\ge1}
   \frac{z_1^{n_1}\cdots z_\ell^{n_\ell}}
        {n_1^{p_1}\cdots n_\ell^{p_\ell}}.
\end{equation}
Its depth is $\ell$. In the disk, all identities below use the
common power-series germ; rational Laurent prefactors have removable
singularities only after the full expression is assembled.

\begin{theorem}[Harmonic rank controls cutoff-generator depth]
\label{thm:depth}
Assume every $a_\nu$ is positive rational, and choose a common denominator
$D$ for the pole locations. If the harmonic pencil has rank $\rho$, then
\begin{equation}\label{eq:generator-general}
 G_R(z)=\sum_{N\ge0}S_{r,N}(R)z^{N+1},\qquad |z|<1,
\end{equation}
can be expressed by finitely many cyclotomic multiple polylogarithms
of level dividing $D$ and depth at most $\rho+1$, rational functions
of $z^{1/D}$, and Euler derivatives of such expressions.
Every positive Euler primitive, normalized to vanish at $z=0$, obeys
the same depth bound.

For integer-grid poles, no fractional power or root-of-unity filter
is needed. In particular, the rank-one families in
Theorems~\ref{thm:classification}--\ref{thm:critical} have depth at
most two. A zero pencil has depth at most one.
\end{theorem}
\begin{proof}
Proposition~\ref{prop:degree} reduces each coefficient to a finite sum
of terms
\begin{equation}\label{eq:coefficient-type}
 Q(N)\prod_{j=1}^{\ell}H_N^{(q_j)}(a_{\nu_j}),\qquad \ell\le\rho,
\end{equation}
where $Q$ is rational and has rational-grid poles. All such sequences
have at most polynomial-logarithmic growth, so their generating series
converge normally in the disk.

First consider integral shifts. Expand the product of finite harmonic
sums by partitioning the index set according to equality and strict
ordering. This is the finite quasi-shuffle argument; it yields nested
sums of depth at most $\ell$. At an equal index, different shifted
rational factors are reduced by ordinary partial fractions. Integer
shifts of a cutoff can be removed by subtracting a finite boundary strip.
Removing one such strip either changes a rational coefficient or reduces
the number of nested free indices; it cannot increase depth.

Next decompose $Q(N)$ into a polynomial and partial fractions. For an
unshifted denominator, the outer summation adds at most one strict
index to the nested harmonic sum. Its diagonal contribution has smaller
depth. For example,
\begin{align}
 \sum_{M\ge1}\frac{z^M H_M^{(q)}}{M^p}
    &=\Li_{p,q}(z,1)+\Li_{p+q}(z),\label{eq:basic-nesting}\\
 \sum_{M\ge1}z^M H_M^{(q)}&=\frac{\Li_q(z)}{1-z}.
 \label{eq:basic-harmonic-gf}
\end{align}
Shifting an outer denominator replaces $M$ by $M+j$, followed by the
same finite boundary-strip procedure. Polynomial factors in $M$ are
implemented by the Euler operator $z\partial_z$. These operations
prove the bound $\ell+1$ on an integer grid.

For rational shifts write $a=A/D$ after moving the representative
into $0<a\le1$ and accounting for finitely many initial terms.
Then
\[
 H_N^{(q)}(A/D)=D^q
    \sum_{\substack{1\le m\le DN+A\\m\equiv A\ (D)}}m^{-q}.
\]
Resolve a congruence by the exact filter
\[
 \mathbf1_{m\equiv A\ (D)}
     =\frac1D\sum_{j=0}^{D-1}\omega^{j(m-A)},
 \qquad \omega=e^{2\pi i/D}.
\]
After rescaling indices, strict inequalities between the original
integer indices differ from strict inequalities between the new
residue-class indices only on equal-block diagonals. Those diagonals
are reduced by partial fractions and have fewer free indices. Induction
on their number reduces to strictly nested sums of the form
\eqref{eq:MPL}, with roots of unity in the arguments and $z^{1/D}$
in the outer argument. Depth does not increase. This also explains
why different residue classes can be handled together.

Finally, a positive Euler primitive divides the coefficient of
$z^{N+1}$ by a positive power of $N+1$. It changes only the rational
factor $Q(N)$ in \eqref{eq:coefficient-type}, leaving the maximum
number of harmonic factors unchanged. The same reduction applies.
\end{proof}

This is a representation theorem, not an irreducibility theorem. It
also does not assert that every endpoint reduces to ordinary zeta values
and their products. Repeated poles introduce inner indices larger than
one, for which the endpoint reduction problem is different from the
simple-pole family in the earlier report.

\subsection{An explicit repeated-pole \texorpdfstring{$\Li_{2,2}$}{Li(2,2)} identity}
Take $r=2$, $a=1$, $s=2$, and write
\[
 S_N=\Z_{2,N}(2,1),\qquad M=N+1.
\]
The weight has repeated poles,
\[
 w_n(2,1)=\frac12\left\{\frac1{(n+1)^2}
                    -\frac2{(n+2)^2}+\frac1{(n+3)^2}\right\}.
\]
A finite calculation from \eqref{eq:finite-moment} gives
\begin{equation}\label{eq:cutoff-PQ}
 S_{M-1}=P(M)+Q(M)H_M^{(2)},\qquad M\ge1,
\end{equation}
where
\begin{align}
 P(M)={}&-\frac7{16}-\frac3{4(M+2)}+\frac7{4(M+2)^2}
                   +\frac7{8(M+1)^2}-\frac1{2(M+1)^4},\label{eq:Pexplicit}\\
 Q(M)={}&\frac38-\frac1{2(M+1)^2}+\frac1{2(M+2)^2}.
 \label{eq:Qexplicit}
\end{align}
For example, $M=1$ gives zero, as required because there are not yet
two distinct indices.

\begin{theorem}[Closed repeated-pole cutoff generator]\label{thm:explicitGF}
For $|z|<1$,
\begin{align}
 \sum_{M\ge1}S_{M-1}z^M
 ={}&\frac{1-z}{2z^2}\Li_{2,2}(z,1)
  +\frac{10-7z}{8z^2(1-z)}\Li_2(z)
  +\frac{4z-7}{4z^2}\Li_1(z)\notag\\
 &-\frac{\Li_4(z)}{2z}+\frac1{2z}-\frac7{16(1-z)}.
 \label{eq:explicitGF}
\end{align}
The apparent singularity at $z=0$ is removable, and the resulting
function begins at order $z^2$.
\end{theorem}
\begin{proof}
For completeness, the three finite moments, with $H=H_M^{(2)}$, are
\begin{align*}
 \mu_{0,M-1}&=\frac{M(3M^3+18M^2+39M+28)}{8(M+1)^2(M+2)^2},\\
 \mu_{1,M-1}&=\frac{M(M-1)(M^2+7M+8)}{8(M+1)^2(M+2)^2},\\
 \mu_{2,M-1}&=H-
  \frac{9M^4+70M^3+133M^2+76M}{8(M+1)^2(M+2)^2}.
\end{align*}
The determinant $\mu_0\mu_2-\mu_1^2$ gives
\eqref{eq:cutoff-PQ}--\eqref{eq:Qexplicit} by rational arithmetic.

The shifted harmonic sums needed to sum that expression are
\begin{align*}
 \sum_{M\ge1}\frac{z^M H_M^{(2)}}{(M+1)^2}
     &=z^{-1}\Li_{2,2}(z,1),\\
 \sum_{M\ge1}\frac{z^M H_M^{(2)}}{(M+2)^2}
     &=z^{-2}\left\{\Li_{2,2}(z,1)
          -(1+z)\Li_2(z)-2(1-z)\Li_1(z)+3z\right\}.
\end{align*}
To verify the second line, put $L=M+2$ and remove the missing term
from $H_{L-1}^{(2)}$. The remaining rational factor is reduced using
\[
 \frac1{L^2(L-1)^2}
 =\frac1{L^2}+\frac1{(L-1)^2}+\frac2L-\frac2{L-1}.
\]
The plain shifted sums are
\[
 \sum_{M\ge1}\frac{z^M}{(M+j)^p}
  =z^{-j}\left\{\Li_p(z)-\sum_{\ell=1}^{j}\frac{z^\ell}{\ell^p}\right\}.
\]
Insert these identities and \eqref{eq:basic-harmonic-gf} into
\eqref{eq:cutoff-PQ}. Collecting terms proves \eqref{eq:explicitGF}.
All steps are absolutely convergent in the disk. The original series
proves removability and its initial vanishing; no separate branch
assignments at the apparent poles are made.
\end{proof}

As $z\uparrow1$, Abel's theorem and the convergent sequence $S_N$ give
\begin{equation}\label{eq:AbelGF}
 \lim_{z\uparrow1}(1-z)\sum_{M\ge1}S_{M-1}z^M
     =\frac38\zeta(2)-\frac7{16}=\frac{\pi^2-7}{16}.
\end{equation}
This independently visible endpoint agrees with
\eqref{eq:concrete-harmonic}. The depth-two coordinate in the cutoff
generator disappears from this particular Abel limit because its
coefficient has an additional factor $1-z$. This is an exact cancellation,
not a general theorem of depth-two endpoint reducibility.

\subsection{Constructing new examples without a search}
Choose any proper rational seed $F$ with poles on positive rational
translated grids and set $R=\Delta^{2r-2}F/(2r-2)!$.
Lemma~\ref{lem:corner-divisibility} proves that its cutoff partition has
an affine harmonic expression and hence a depth-two generator at most.
This works for seeds with several pole orders and several residue classes.
If the seed is itself a first difference, the harmonic pencil vanishes
and the generator has depth at most one.

For example, any seed of the form
\[
 F(x)=\frac{u}{x+1/3}+\frac{v}{(x+1/2)^2}
                         +\frac{w}{(x+2/3)^3}
\]
produces an all-particle rank-at-most-one family on level six. The corner
entries are $u,v,w$ in their respective harmonic coordinates by
\eqref{eq:lambda-seed}. When $u\ne0$, the decay is nonzero critical
and Theorem~\ref{thm:critical} applies. The theorem supplies a finite
construction and a complete rank certificate; it does not require a
numerical integer-relation search.
